
% !TEX root = main.tex

\usepackage[utf8]{inputenc}

%----------------------------------------------------------------------------------------
%	LOGIC
%----------------------------------------------------------------------------------------

% provides \isempty test
\usepackage{xifthen}

%----------------------------------------------------------------------------------------
%	FONT
%----------------------------------------------------------------------------------------

% some tex-live fonts
% \usepackage[defaultsans]{droidsans}
% \usepackage[default]{comfortaa}
% \usepackage{cmbright}
% \usepackage[default]{raleway}
% \usepackage{fetamont}
% \usepackage[default]{gillius}
%\usepackage[light,math]{iwona}
% \usepackage[thin]{roboto}

% set font default
\renewcommand*\familydefault{\sfdefault}
\usepackage[T1]{fontenc}

% more font size definitions
\usepackage{moresize}

\usepackage{fontawesome}

%----------------------------------------------------------------------------------------
%	PAGE LAYOUT  DEFINITIONS
%----------------------------------------------------------------------------------------

% debug page outer frames
%\usepackage{showframe}


% define page styles using geometry
\usepackage[a4paper]{geometry}
\geometry{top=1cm, bottom=-.6cm, left=0.4cm, right=1cm}


%less space between header and content
\setlength{\headheight}{-5pt}


%indentation is zero
\setlength{\parindent}{0mm}

%----------------------------------------------------------------------------------------
%	TABLE /ARRAY DEFINITIONS
%----------------------------------------------------------------------------------------

%for layouting tables
\usepackage{multicol}
\usepackage{multirow}

%extended aligning of tabular cells
\usepackage{array}

\newcolumntype{x}[1]{
    >{\raggedleft\hspace{0pt}}p{#1}
}


%----------------------------------------------------------------------------------------
%	GRAPHICS DEFINITIONS
%----------------------------------------------------------------------------------------

\usepackage{graphicx}
\graphicspath{
    {imgs/}
}

\usepackage{tcolorbox}
%for floating figures
\usepackage{wrapfig}
\usepackage{float}
%\floatstyle{boxed}
%\restylefloat{figure}

\usepackage{tikz}
\usetikzlibrary{shapes, backgrounds, mindmap, trees}

%----------------------------------------------------------------------------------------
%	Color DEFINITIONS
%----------------------------------------------------------------------------------------
\usepackage{transparent}
\usepackage{color}

%accent color
\definecolor{complcol}{RGB}{250,150,10}

%dark background color
\definecolor{bgcol}{RGB}{110,110,110}

%light background / accent color
\definecolor{softcol}{RGB}{225,225,225}

\definecolor{sectcol}{RGB}{0,120,150}


%Package for links, must be the last package used
\usepackage[hidelinks]{hyperref}



%============================================================================%
%
%
%	DEFINITIONS
%
%
%============================================================================%

% returns minipage width minus two times \fboxsep
% to keep padding included in width calculations
\newcommand{\mpwidth}{\linewidth-\fboxsep-\fboxsep}

%----------------------------------------------------------------------------------------
% CUSTOM STRUT FOR EMPTY BOXES
%----------------------------------------- -----------------------------------------------
\newcommand{\mystrut}{
    \rule[-.3\baselineskip]{0pt}{\baselineskip}
}

%----------------------------------------------------------------------------------------
% CUSTOM LOREM IPSUM
%----------------------------------------------------------------------------------------
\newcommand{\lorem}{
    Lorem ipsum dolor sit amet, consectetur adipiscing elit. Donec a diam lectus.
}
%----------------------------------------------------------------------------------------
% 	ARROW GRAPHICS in Tikz
%----------------------------------------------------------------------------------------
% a six pointed arrow pointing to the left/right
\newcommand{\tzlarrow}{
    (0,0) -- (0.2,0) -- (0.3,0.2) -- (0.2,0.4) -- (0,0.4) -- (0.1,0.2) -- cycle;
}
\newcommand{\tzrarrow}{
    (0,0.2) -- (0.1,0) -- (0.3,0) -- (0.2,0.2) -- (0.3,0.4) -- (0.1,0.4) -- cycle;
}

% include the left/right arrow into a tikz picture
% param1: fill color
\newcommand{\larrow}[1]{
    \begin{tikzpicture}[scale=0.58]
        \filldraw[fill=#1!100,draw=#1!100!black]  \tzlarrow
    \end{tikzpicture}
}
\newcommand{\rarrow}{
    \begin{tikzpicture}[scale=0.7]
    	\filldraw[fill=sectcol!100,draw=sectcol!100!black] \tzrarrow
    \end{tikzpicture}
}

%----------------------------------------------------------------------------------------
%	Custom sections
%----------------------------------------------------------------------------------------

% create a coloured box with arrow and title as cv section headline
% param 1: section title
\newcommand{\cvsection}[1]{
    \colorbox{sectcol}{\mystrut \makebox[1\mpwidth][l]{
    \larrow{bgcol} \hspace{-8pt} \larrow{bgcol} \hspace{-8pt} \larrow{bgcol} \textbf{\textcolor{white}{\uppercase{#1}}}\hspace{4pt}
    }}\\
}

% create a coloured arrow with title as cv meta section section
% param 1: meta section title
\newenvironment{metasection}[1]{
	\vspace{6pt}
	\begin{center}
		\textcolor{white}{\large{\uppercase{#1}}}\\
	\normalsize
	\parbox{0.7\mpwidth}{\textcolor{white}	\hrule}
}{\end{center}}


% Content will be aligned on the left
\newenvironment{metaSectionLeft}[1]{
	\vspace{6pt}
    \begin{center}
        \textcolor{white}{\large{\uppercase{#1}}}\\
        \begin{minipage}[t]{0.85\linewidth}
            \normalsize
            \parbox{\mpwidth}{\textcolor{white}	\hrule}

}{
        \end{minipage}
    \end{center}
}

% Content will be centred
\newenvironment{metaSectionCenter}[1]{
	\vspace{6pt}
    \begin{center}
        \textcolor{white}{\large{\uppercase{#1}}}\\
        \begin{minipage}[t]{0.85\linewidth}
            \normalsize
            \parbox{\mpwidth}{\textcolor{white}	\hrule}
            \begin{center}
}{
            \end{center}
        \end{minipage}
    \end{center}
}


%----------------------------------------------------------------------------------------
%	 CV EVENT
%----------------------------------------------------------------------------------------

% creates a stretched box as cv entry headline followed by two paragraphs about
% the work you did
% param 1:	event time i.e. 2014 or 2011-2014 etc.
% param 2:	event name (what did you do?)
% param 3:	institution (where did you work / study)
% param 4:	what was your position
% param 5:	some words about your contributions
\newcommand{\cvevent}[5]{
    % \vspace{8pt}
	\begin{tabular*}{1\mpwidth}{p{0.55\mpwidth}  x{0.42\mpwidth}}
        \textcolor{black}{\textbf{#2}} & \textcolor{complcol}{#3}, \textcolor{bgcol}{#1}
	\end{tabular*}
    % \vspace{-12pt}
    \textcolor{softcol}{\hrule}
    \vspace{6pt}
	\begin{tabular*}{0.5\mpwidth}{p{\mpwidth}}
        \larrow{softcol}  #4\\[3pt]
        \larrow{softcol}  #5\\[6pt]
	\end{tabular*}
}

% creates a stretched box as
\newcommand{\cveventmeta}[2]{
	\mbox{\mystrut \hspace{87pt}\textit{#1}}\\
	#2
}

% use to vertically center content
% credits to: http://tex.stackexchange.com/questions/7219/how-to-vertically-center-two-images-next-to-each-other
\newcommand{\vcenteredinclude}[1]{
    \begingroup
    \setbox0=\hbox{\includegraphics{#1}}%
    \parbox{\wd0}{\box0}
    \endgroup
}

% use to vertically center content
% credits to: http://tex.stackexchange.com/questions/7219/how-to-vertically-center-two-images-next-to-each-other
\newcommand*{\vcenteredhbox}[1]{
    \begingroup
    \setbox0=\hbox{#1}\parbox{\wd0}{\box0}
    \endgroup
}

%----------------------------------------------------------------------------------------
%	ICON-SET EMBEDDING
%----------------------------------------------------------------------------------------
%icon shortcut
\newcommand{\icon}[3]{\makebox(#2, #2){\textcolor{#3}{\csname fa#1\endcsname}}}

%icon with text shortcut
\newcommand{\icontext}[4]{
	\vcenteredhbox{\icon{#1}{#2}{#4}}
    \vcenteredhbox{\textcolor{#4}{#3}}
}

%icon with website url
\newcommand{\iconhref}[5]{
    \vcenteredhbox{\icon{#1}{#2}{#5}}
    \href{#4}{\textcolor{#5}{#3}}
}

%icon with email link
\newcommand{\iconemail}[5]{
    \vcenteredhbox{\icon{#1}{#2}{#5}}
    \href{mailto:#4}{\textcolor{#5}{#3}}
}
